\section{Baseline Model}

\paragraph{Preferences}

Following \textcite{Keane2016}, households derive utility from consumption and disutility from labour supply. Period utility is given by

\[
U_t =
\frac{c_t^{1+\eta}}{1+\eta}
-
\beta(n_t)\frac{h_t^{1+\gamma}}{1+\gamma},
\]

where $c_t$ denotes consumption and $h_t$ hours worked. $\eta$ and $\gamma$ determine the curvature of utility from consumption and the disutility of labour. $\beta(n_t)$ scales the disutility of labour depending if children are present.

\paragraph{Recursive formulation}

The household is described in each period $t$ by the state vector
\[
x_t = (n_t, a_t, k_t),
\]
where $n_t\in\{0,1\}$ is the number of children, $a_t$ assets, and $k_t$ human capital. The controls are consumption $c_t$ and hours $h_t$. Taking the expectation $\mathbb{E}_t$ over next period's fertility shock, the value function solves the Bellman equation

\[
V_t(n_t,a_t,k_t)
=
\max_{c_t,h_t}
\left\{
\frac{c_t^{1+\eta}}{1+\eta}
-
\beta(n_t)\frac{h_t^{1+\gamma}}{1+\gamma}
+
\rho \, \mathbb{E}_t
\left[
V_{t+1}(n_{t+1},a_{t+1},k_{t+1})
\right]
\right\},
\]

where $\rho$ is the discount factor.

\paragraph{State transitions}

The three states evolve as follows. Fertility follows a stochastic process,

\[
n_{t+1} =
\begin{cases}
n_t + 1 & \text{with probability } p(n_t), \\
n_t & \text{with probability } 1 - p(n_t),
\end{cases}
\qquad
p(n_t)=
\begin{cases}
p_n & \text{if } n_t = 0, \\
0 & \text{if } n_t = 1,
\end{cases}
\]
so at most one child can arrive. Assets evolve according to the budget constraint

\[
a_{t+1} = (1+r)a_t + (1-\tau) w_t h_t - c_t,
\]

where $\tau$ is the labour income tax rate and the after-tax wage is $w_t = w(1+\alpha k_t)$. Human capital accumulates through labour supply,

\[
k_{t+1} = k_t + h_t,
\]

and the disutility of labour depends on children through

\[\beta(n_t)=\beta_0+\beta_1 n_t.\]

We solve by backwards induction, assuming no bequest motive. With no continuation value, the last-period problem reduces to a static choice of hours,

\[V_{T}(n_T,a_T,k_T)=\max_{h_{T}}\frac{c_{T}^{1+\eta}}{1+\eta}-\beta(n_T)\frac{h_{T}^{1+\gamma}}{1+\gamma} \quad \text{s.t} \quad 
c_T = a_{T}+(1-\tau)w_{T}h_{T}.
\]